\section{Discussion and Limitations}
\label{sec:discussion}
\paragraph{Formal and native scope.}
The formal result separates conditional recovery from single-challenge
index privacy. The latter assumes distribution-specific uniform-secret
PLWE hardness and efficient sampling; it neither follows from empirical
correctness nor requires a successful decryption event. The rounded
Gaussian width remains symbolic. Its exact efficient machine realization,
approximation and computational randomness require independent evidence.
The native bounded, biased sampler and OS-seeded PRNG have not been
bridged to that law, and a public native noise parameter is not a proved
absolute bound on integer coefficients. Consequently neither finite
native passes nor the parametric bound supplies a native negligible-failure
claim or a contemporary concrete-security certificate.

After a suitable joint-law and implementation bridge, task accounting
could take the form
\(\epsilon_{\rm task}\le\epsilon_{\rm dec}+d_{\rm joint}
+\epsilon_{\rm other}\), with separately justified sampling, abort and
runtime terms. An abort included in the joint-law distance is not counted
again. PRNG replacement needs a separate computational failure-test bound.
No such bound is established here. Independent repeated requests must
also share one complete-task reliability budget, as detailed in the
supplement. The privacy game gives neither symmetric PIR nor protection
against malicious replies, adaptive key reuse or decryption oracles.

\paragraph{Empirical generality.}
The measurements cover one flat backend path, one WSL2 host and finite
workloads. The comparator uses concurrent private workers under the same
profile within each experiment; it is not a globally optimized XPIR
service with shared persistent preprocessing. External PIR timings are
not aligned measurements. The capacity study adds one two-block length,
and bundle sensitivity covers only one workload. Neither supports
arbitrary-block, cross-host or bare-metal extrapolation.

The primary and multiplicity resource limits were enforced by their
runners, but individual monitoring samples were not retained. Their
transient virtual-memory and swap behavior cannot be reconstructed from
the surviving endpoints and summaries. The sensitivity/capacity data
retain actual monitor timestamps, yet sampled maxima are still not
instantaneous peaks. This difference limits resource diagnosis without
establishing invalidity of the recorded task times. Bootstrap intervals
are pointwise; three successive sessions give descriptive variation.
The unexplained two-record effect difference remains part of the evidence.

\paragraph{Design tradeoff.}
Radix packing places several bounded values in each coefficient without
bucket placement or additional SIMD rotations, while consuming plaintext
capacity and no-wrap margin. It retains a linear-size query and does not
claim general communication efficiency. A deployment would require a
joint choice of correctness and security parameters, a justified sampler
and randomness implementation, transport measurements, and comparisons
with optimized preprocessing and query-compression alternatives.
